\documentclass[11pt,a4paper]{article}
\usepackage[utf8]{inputenc}
\usepackage[T1]{fontenc}
\usepackage{amsmath, amssymb}
\usepackage{graphicx}
\usepackage{tikz}
\usepackage[margin=2.5cm]{geometry}
\usepackage{hyperref}

\title{The Title of Your Article}
\author{Your Name \\ \small Your Institution}
\date{\today}

\begin{document}
\maketitle

\begin{abstract}
One paragraph that says what the article is about, what was done, and what was found.
\end{abstract}

\section{Introduction}
\label{sec:intro}
Start here. Cite earlier work like this \cite{knuth1984}, and point to a later
section like this: Section~\ref{sec:results}.

\section{Method}
An equation, numbered so it can be referred to:
\begin{equation}
  \label{eq:energy}
  E = m c^2 .
\end{equation}
Equation~\eqref{eq:energy} is referred to by its label.

\section{Results}
\label{sec:results}
Figure~\ref{fig:sketch} is drawn with TikZ, so it needs no image file.

\begin{figure}[h]
  \centering
  \begin{tikzpicture}
    \draw[->] (0,0) -- (4.2,0) node[right] {$x$};
    \draw[->] (0,0) -- (0,3.2) node[above] {$y$};
    \draw[thick, blue] (0,0.3) .. controls (1.5,3) and (2.5,0.5) .. (4,2.6);
  \end{tikzpicture}
  \caption{A sketch of the result.}
  \label{fig:sketch}
\end{figure}

\section{Conclusion}
What the reader should take away.

\bibliographystyle{plain}
\bibliography{references}

\end{document}
